\chapter{The Skeletal and Muscular Systems}
\section{Get Rolling}
\textbf{Name the organs and then match them with the organ system they are part of.}
\begin{enumerate}
    \item heart - circulatory system
    \item kidney - excretory system
    \item skull - skeletal system
    \item liver - digestive system 
\end{enumerate}
\section{\textcolor{red}{What to Know}}
Different parts of the skeletal system, different types of joints and how muscles make movements possible
\section{\textbf{Learning Objectives}}
\begin{itemize}
    \item to define the skeletal system
    \item to list parts of the skeleton
    \item to distinguish between different types of joints
    \item to identify some cartilages in our body
    \item to describe the working of muscles
\end{itemize}
\section{I Explore}
Our body has different \textbf{organ systems.} One such organ system is the skeletal system.
\vspace{0.1in}
\par
Before moving on to the skeletal system, let us learn about the major organ that constitutes the skeletal system, that is, the bone.
\section{\textcolor{red}{Bones}}
An adult human body has 206 bones, whereas a newborn has around 300 bones. As a child grows into an adult, many bones \textcolor{red}{fuse} together. The bones vary in their shapes and sizes. For example, bones of our head are flat and round, whereas bones of our legs and arms are long.
\vspace{0.1in}
\par
The \textcolor{red}{framework} of bones in a human body constitutes the \textbf{skeleton.}
\vspace{0.1in}
\par
The skeleton gibes shape to our body.
\vspace{0.1in}
\par
The \textbf{skeletal system} comprises the skeleton, cartilage and other connective tissues that support the body and give it shape.
\vspace{0.3in}
\textcolor{red}{I wonder} what would happen if I had no skeleton!
\section{\textcolor{red}{Skeleton and its Parts}}
The human skeleton consists of the skull, backbone, ribcage and two pairs of limbs-forelimbs or arms and hindlimbs or legs.
\subsection{\textcolor{orange}{Skull}}
The skull is the bony structure that protects our brain. The skull is made up of hard and tough bones connected together. There are 22 bones in our skull-8 bones on the upper part of the skull and 14 bones on the face.
\vspace{0.1in}
\par
All the bones of the skull except the lower jawbone are fixed. The movable lower jawbone enables us to eat and speak.
\subsection{\textcolor{orange}{Backbone}}
The backbone is also known as the vertebral column. You can feel your backbone by touching the center of your back and moving your hand up and down.
\vspace{0.1in}
\par
It is a long, \textcolor{red}{flexible} column made up of 33 individual bones called \textbf{vertebrae.}
\vspace{0.1in}
\par
These bones help us bend, twist and turn.
\vspace{0.1in}
\par
Each vertebra is a small, round bone that is connected to the ones above and below it.
\vspace{0.1in}
\par
These vertebrae are hollow at the center, which forms a long tunnel through which a vital part of our body, called the \textbf{spinal cord,} runs.
\subsection{\textcolor{orange}{Ribcage}}
The ribcage is a bony structure located in our chest area. It has curved bones called the \textbf{ribs} that form a cage-like structure. in our chest. It protects our chest organs, such as the heart and the lungs. We have 12 pairs of ribs in total. The ribs are connected to our backbone at the back and to the breastbone (sternum) at the front. The last two pairs are not attached to the sternum and thus are called the \textbf{floating ribs.}
\subsection{\textcolor{orange}{Limbs}}
The bones in our limbs help us move and perform various activities. In humans, there are two pairs of limbs-forelimbs (arms) and hindlimbs (legs).
\subsection{\textbf{Forelimbs}}
A \textbf{forelimb} consists upper arm and a lower arm. These two parts of an arm meet at the elbow.
\vspace{0.1in}
\par
The upper arm has one bone called the \textbf{humerus} and the lower arm has two bones called \textbf{radius} and \textbf{ulna.} The wrist and the hand are made of many small bones.
\vspace{0.9in}
\subsection{\textbf{Hindlimbs}}
The \textbf{hindlimbs} also have three bones. The thigh bone is the longest bone in the body. It is also called the \textbf{femur,} the longest bone in the body. The femur is attached to the hip bone. The lower leg has two bones called the \textbf{tibia} and the \textbf{fibula.} The ankle and toes are made of many small bones.
\subsection{\textcolor{orange}{Cartilage}}
Try to fold your ears. See if you can do so. Yes, you can. It is because of the cartilage
\vspace{0.1in}
\par
\textbf{Cartilage} is a smooth and flexible tissue that covers the ends of bones at joints. It acts as a cushion and reduces friction between bones during movement. It is softer and much more flexible than a bone.
\vspace{0.3in}
\textcolor{red}{I wonder} what would happen if our body did not have cartilages!
\vspace{3.1in}
\section{\textcolor{red}{Joints and their Types}}
\textbf{Joints} are the places where two or more bones come together in our body. They allow our bones to move and bend, making it possible for us to do all kinds of movements.
\vspace{0.1in}
\par
There are two main types of joints, movable joints and immovable joints. Let us see how they are different from each other.\\
\vspace{0.1in}
\begin{tabular}{|c|c|}
\centering
\hline
Immovable Joints & Movable Joints\\
\hline
Immovable joints, as the name suggests, do not move. & Movable joints are the ones that allow our bodies to move in different ways.\\
\hline
These joints are found in places like our skull and the bones in our ribcage. & These joints are found in places like our arms, legs and fingers.\\
\hline
They are hard and help protect important parts of our body. & They help us bend, twist and move around easily.
\hline
\end{tabular}
\vspace{0.1in}
\par
There are different types of movable hoints in our body. Let us learn about them.
\subsection{\textcolor{orange}{Hinge Joints}}
\textbf{Hinge Joints} work like a door hinge. They allow movement in only one direction, like a door opening and closing. Knees and elbows are examples of hinge joints.
\subsection{\textcolor{orange}{Ball-and-socket joints}}
\textbf{Ball-and-socket joints} are like a ball fitting into a socket. They allow movement in many directions. Our shoulder and hip joints are ball-and-socket joints. They enable us to swing our arms and legs around.
\subsection{\textcolor{orange}{Pivot Joints}}
Pivot Joints let bones rotate or twist. Our neck has a pivot joint, allowing us to turn our head from side to side.
\subsection{\textcolor{orange}{Gliding Joints}}
\textbf{Gliding joints} let bones slide in different directions. We have them in our wrists and ankles, helping us move our hands and feet in various ways.
\vspace{2.1in}
\section{\textcolor{red}{Muscles}}
We have muscles in our bodies that help us move. These muscles are connected to our bones and work together, allowing us to do activities like running, jumping and smiling. Our body has more than 600 muscles.
\vspace{0.1in}
\par
The muscles attached to the bones by strong fibres called \textbf{tendons.}
\vspace{0.1in}
\par
There are majorly two kinds of muscles-voluntary muscles and involuntary muscles. The muscles that are in our control are called \textbf{voluntary muscles} and the ones which are not in our control are called \textbf{involuntary muscles.}
\subsection{\textcolor{orange}{How Do Muscles Work?}}
Let us understand how our muscles work, taking the example of our arm. The arms have two important muscles called the biceps and the triceps.
\vspace{0.1in}
\par
The \textbf{biceps} are on the front side of the upper arm. Try bending your arm. When your biceps contract, they become active and pull the forearm up towards the shoulder.
\vspace{0.1in}
\par
The \textbf{triceps} are on the back side of the upper arm. They help us straighten the arm at the elbow.
\vspace{0.1in}
\par
Let us take the example of lifting a weight. When we lift a weight, the biceps contract and the triceps relax. When we put down the weight, the triceps contract and the biceps relax.
\vspace{1.1in}
\subsection{\textcolor{orange}{Care for Bones and Muscles}}
To maintain strong bones and muscles, we should:
\vspace{0.1in}
\begin{itemize}
    \item exercise every day.
    \item play in sunlight (source of vitamin D).
    \item eat food rich in calcium.
    \item maintain good posture.
\end{itemize}
\section{Now I Know}
\begin{itemize}
    \item A group of organs that works together to perform a function forms an \textbf{organ system.}
    \item An adult human body has 206 bones, whereas a newborn has around 300 bones.
    \item The framework of bones in a human body constitutes the \textbf{skeleton.}
    \item The \textbf{skeletal system} comprises the skeleton, cartilage and other connective tissues.
    \item The human skeleton consists of the skull, the backbone, the ribcage and two pairs of limbs-forelimbs or arms and hindlimbs or legs.
    \item \textbf{Joints} are the places where two or more bones come together in our body. They allow our bones to move and bend, making it possible for us to do all kinds of movements.
    \item Some of the movable joints include hinge joints, ball-and-socket joints, pivot joints and gliding joints.
    \item Muscles are attached to our bones by strong fibres called \textbf{tendons,} and they work together to help us move in various ways.
    \item There are majorly two kinds of muscles-voluntary muscles and involuntary muscles.
\end{itemize}